
\documentclass[11pt]{amsart}

\usepackage[a4paper,margin=1in]{geometry}
\usepackage{amsmath,amssymb,amsthm,mathtools}
\usepackage{bm}
\usepackage{booktabs}
\usepackage{enumitem}
\usepackage[colorlinks=true,linkcolor=blue,citecolor=blue,urlcolor=blue]{hyperref}

\newtheorem{theorem}{Theorem}
\newtheorem{proposition}{Proposition}
\newtheorem{lemma}{Lemma}
\newtheorem{corollary}{Corollary}
\newtheorem{conjecture}{Conjecture}
\newtheorem{definition}{Definition}
\newtheorem{remark}{Remark}
\newtheorem{openproblem}{Open Problem}

\newcommand{\Rm}{\mathrm{Rm}}
\newcommand{\Ric}{\mathrm{Ric}}
\newcommand{\Ricci}{\mathrm{Ric}}
\newcommand{\R}{\mathrm{R}}
\newcommand{\tr}{\operatorname{tr}}
\newcommand{\grad}{\operatorname{grad}}
\newcommand{\dmu}{\,d\mu_g}
\newcommand{\dt}{\partial_\tau}
\newcommand{\lP}{\ell_P}
\newcommand{\KP}{K_{\!P}}
\newcommand{\F}{\mathcal F}
\newcommand{\B}{\mathcal B}
\newcommand{\U}{\mathcal U}
\newcommand{\Wcal}{\mathcal W}
\newcommand{\Ncal}{\mathcal N}
\newcommand{\Mcal}{\mathcal M}
\newcommand{\Vol}{\operatorname{Vol}}
\newcommand{\inj}{\operatorname{inj}}
\newcommand{\Grad}{\operatorname{Grad}}



\title[Supercritical Variational Curvature Barriers]{Supercritical Variational Curvature Barriers for Ricci-Type Flows on Compact \(S^3\)}
\author{B.~Batenin}
\address{Independent researcher, Berlin, Germany}
\author{A.~Preece}
\address{Independent researcher, Berlin, Germany}
\date{}
\subjclass[2020]{53E20, 53C44, 58J35}
\keywords{Ricci-type flow, curvature barrier, Berger spheres, Perelman entropy, noncollapse}

\begin{document}
\raggedbottom

\begin{abstract}

This paper is not a revision of the operational Planck-boundary interpretation
used in the compact \(S^3\) programme.  In that setting, \(\ell_P\) marks the
limit of physical distinguishability, not a dynamical forcing term in the Ricci
flow.  The present note asks a separate mathematical question: if Ricci-type
flows are treated as autonomous geometric dynamics on compact \(S^3\), which
variational curvature barriers can suppress supercritical curvature growth
before a prescribed curvature scale is reached?

A Hopf--Berger analysis of Ricci-type flows on compact \(S^3\)
shows that bounded traceless forcing, fibre-restoring forcing, and a constant
cosmological-constant term do not supply the negative cubic damping needed to
control the round-collapse direction.  The homogeneous term
\(h=2\lambda |\Rm|g\) has the desired sign in a finite-dimensional reduction,
but fails as a PDE mechanism because the discarded Laplacian source
\(-4\lambda \Delta |\Rm|\) is of the same scale as the intended damping near a
curvature concentration.

This paper isolates the corresponding variational replacement.  The local
high-curvature regulator is the negative \(L^2\)-gradient of a curvature
energy
\[
        \F_\Phi[g]=\int_{S^3}\Phi(K)\,d\mu_g,\qquad K=|\Rm|^2,
\]
and, in the singular wall version,
\[
        \B_{p,q}[g]
        =
        \int_{S^3}
        \frac{(K+\delta)^q}{(\KP-K)^p}\,d\mu_g,
        \qquad
        \KP=\ell_P^{-4}.
\]
The algebraic high-curvature test is passed.  In dimension three
\[
        K=|\Rm|^2=4|\Ric_0|^2+\frac13R^2,
\]
so \(K\) detects both round and anisotropic curvature, and the conformal
variation gives the decisive factor
\[
        4K\Phi'(K)-3\Phi(K).
\]
For \(\Phi(K)=K^q\) this is \((4q-3)K^q\), so the sharp scale threshold in
dimension three is \(q>3/4\).  For the Planck wall the same factor diverges
with the correct sign as \(K\nearrow\KP\).  On the full Berger family the wall
has an exact two-dimensional reduction and the anisotropy ratio \(u=z/x\)
satisfies
\[
        \dot u
        =
        \frac{8u(1-u)}{x}
        \left(1+\frac{2\varepsilon\Phi'(K)}{x}(11u-4)\right),
\]
which yields an invariant cone \(u\ge4/11\) and identifies the lower Hopf
sector where the wall may accelerate fibre collapse.  The same reduction gives
closed formulas for volume, scalar curvature, traceless Ricci size, and the
signs of the Planck-wall Euler tensor on Berger sectors.  A simple zero-mode
heat-kernel lower bound also shows that the scale detector
\(\sup_{0<t\le r^2}t^{3/2}H_g(t,x,x)\) diverges along fixed-base Hopf
collapse.  A wall gap gives an explicit Ricci lower bound and hence controlled
heat-kernel comparison constants, although no propagation theorem for this
detector is claimed.

The paper also records the exact obstruction to turning this wall into a
noncollapse theorem by an additive Perelman-potential repair.  On the
homogeneous Berger sector the candidate
\(\lambda_\Phi=\lambda_1(-4\Delta+R+\varepsilon\Phi(K))\) reduces exactly to
\(R+\varepsilon\Phi(K)\).  For \(\Phi(K)=K\), at the round point,
\[
        \dot\lambda_\Phi
        =
        \frac{12(2x-\varepsilon)(x+4\varepsilon)}{x^4},
\]
so monotonicity fails for every fixed \(\varepsilon>0\) once
\(x<\varepsilon/2\).  More generally, additive matching \(K\)-potential repairs
of the form
\[
        \Wcal_{\mathrm{Perelman}}+\sigma\varepsilon\tau^\alpha K
\]
fail the same Berger test for every sign \(\sigma\) and every exponent
\(\alpha\).  For the singular Planck wall
\(\Phi(K)=(K+\delta)^q(\KP-K)^{-p}\), the corresponding additive repair
also fails analytically at the wall: its Berger derivative has a negative
leading term of order
\[
        -\varepsilon^2(\KP-K)^{-2p-2}.
\]
Thus the curvature wall is a valid local high-curvature metric regulator, but
it is not itself a noncollapse entropy.  By contrast, the unmodified Perelman
entropy passes the protected positive-scalar Berger cone test: on
\(4/11\le u<4\) its reduced derivative along the wall flow is nonnegative.
This gives the first positive non-additive two-mechanism result in the paper.
A complete internal replacement for surgery would still require a corresponding
propagation or monotonicity theorem for the unmodified entropy along the full
inhomogeneous DeTurck-gauged wall flow.
\end{abstract}

\maketitle

\section{Purpose}

The status of this paper is deliberately mathematical.  It does not claim that
the operational Planck-boundary interpretation requires a curvature-wall
correction, nor does it replace that interpretation.  In the compact \(S^3\)
Planck-boundary programme, \(\ell_P\) marks a limit of physical
distinguishability.  Here we ask a different question: once modified
Ricci-type flows are studied as autonomous geometric systems, can one construct
intrinsic variational forcing terms that prevent supercritical curvature growth
before a prescribed curvature scale is reached?

The object is a Ricci-type flow on a compact closed three-sphere
\[
        \dt g=-2\Ric+h[g].
\]
The problem is not to imitate Ricci-flow surgery.  Surgery is an external
topological intervention.  The question here is narrower: can an internal,
geometrically natural forcing term prevent a high-curvature collapse inside a
fixed closed slice?

The useful test is the curvature reaction estimate
\[
        (\dt-\Delta)|\Rm|^2
        =
        -2|\nabla\Rm|^2+C_{\mathrm{RF}}(\Rm),
        \qquad
        |C_{\mathrm{RF}}(\Rm)|\le C_3|\Rm|^3 .
\]
A local forcing can dominate the singular regime only if, at large curvature,
its curvature contribution contains a negative term of order strictly stronger
than, or at least cubic with coefficient exceeding, the Ricci reaction:
\[
        2\langle C_h(\Rm),\Rm\rangle
        \le
        -\gamma |\Rm|^3+A|\Rm|^2+B,
        \qquad
        \gamma>C_3.
\]
A constant cosmological term does not pass this test.  If \(h=2\Lambda g\),
then its contribution is
\[
        -4\Lambda |\Rm|^2,
\]
which is quadratic.  It cannot dominate \(C_3|\Rm|^3\) as
\(|\Rm|\to\infty\).

The preceding Hopf--Berger calculation also gives a geometric warning.  For
\[
        g_{a,c}=a^2(\sigma_1^2+\sigma_2^2)+c^2\sigma_3^2
        \quad \text{on } SU(2)\simeq S^3,
\]
one has
\[
        |\Rm|^2
        =
        \frac{4(16a^4-24a^2c^2+11c^4)}{a^8}.
\]
Equivalently, with \(x=a^2\), \(z=c^2\), and \(u=z/x\),
\[
        K:=|\Rm|^2
        =
        \frac{4(16-24u+11u^2)}{x^2}.
\]
Thus \(c\to0\) at fixed \(a\) is a bounded-curvature fibre collapse, while
\(a\to0\) at fixed \(u\) is the genuine curvature singularity.  Any universal
regularizer must see the round/base collapse.  Traceless-only or fibre-only
invariants do not suffice.

\section{Operational meaning of universality}

The word ``universal'' must be made finite.  We use it in the following
operational sense.

\begin{definition}[local high-curvature regularizer]
A forcing \(h[g]\) is called a local high-curvature regularizer on compact
three-manifolds if it satisfies four tests.

\begin{enumerate}[label=(\roman*)]
\item It is natural: \(h[\phi^*g]=\phi^*h[g]\) for diffeomorphisms \(\phi\).
\item It is local: it is built from \(g\), curvature, and finitely many
covariant derivatives.
\item Its scalar detector does not vanish on a non-flat three-dimensional
curvature tensor.
\item In a high-curvature concentration it supplies negative leading-order
curvature drift strong enough to dominate \(C_3|\Rm|^3\), without discarding
spatial derivative terms of the same order.
\end{enumerate}
\end{definition}

This definition deliberately excludes a common false success.  A term such as
\(2f(|\Rm|)g\) may look good in a homogeneous reduction, but its scalar
curvature variation includes Laplacians of \(f(|\Rm|)\).  For
\(f=\lambda |\Rm|\),
\[
        \delta R=-2fR-4\Delta f
\]
in dimension three.  The homogeneous model keeps \(-2\lambda |\Rm|R\) and
drops \(-4\lambda\Delta |\Rm|\).  Near a neck concentration of width \(L\) and
peak \(|\Rm|_{\mathrm{pk}}\sim L^{-2}\), the dropped term scales like
\[
        \lambda \frac{|\Rm|_{\mathrm{pk}}}{L^2}
        \sim
        \lambda |\Rm|_{\mathrm{pk}}^2,
\]
the same order as the intended damping.  It is therefore not a stable PDE
mechanism.

The replacement is to keep the derivative terms by making the forcing
variational.

\section{The detector: \texorpdfstring{\(K=|\Rm|^2\)}{K=|Rm|2}}

In dimension three the Weyl tensor vanishes, and the full curvature tensor is
determined by \(\Ric\).  The norm identity is
\[
        |\Rm|^2
        =
        4|\Ric|^2-R^2
        =
        4|\Ric_0|^2+\frac13R^2 .
\]
Here
\[
        \Ric_0=\Ric-\frac13Rg
\]
is the traceless Ricci tensor.

\begin{lemma}[minimal scalar detector]
In dimension three, \(K=|\Rm|^2\) is the minimal positive scalar curvature
detector for the present problem: it detects both the round direction
\((\Ric_0=0,\ R\ne0)\) and anisotropic curvature
\((\Ric_0\ne0)\).  A traceless-only detector misses the round collapse; scalar
curvature alone does not give a positive definite norm of the whole curvature
tensor; derivative detectors such as \(|\nabla\Rm|^2\) miss homogeneous
round collapse.
\end{lemma}

\begin{proof}
The identity above is pointwise.  The term \(R^2/3\) is precisely the round
component, while \(4|\Ric_0|^2\) is the anisotropic component.  Removing either
term loses one of the two curvature channels.  A derivative-only detector
vanishes on homogeneous Berger metrics, including the round shrinking mode,
even while \(K\to\infty\).
\end{proof}

Thus the basic invariant must be a function of
\[
        K=|\Rm|^2.
\]

\section{Variational forcing}

Let
\[
        \F_\Phi[g]=\int_M \Phi(K)\dmu,
        \qquad K=|\Rm|^2,
\]
where \(\Phi\) is smooth on the relevant interval.  Define \(E_\Phi(g)\) by
\[
        D\F_\Phi[g](v)
        =
        \int_M \langle E_\Phi(g),v\rangle_g\,d\mu_g .
\]
The proposed forced flow is
\[
        \boxed{\dt g=-2\Ric-\varepsilon E_\Phi(g)}
\]
with \(\varepsilon>0\).

This is a fourth-order geometric flow, because \(K\) depends on two
derivatives of the metric and \(E_\Phi\) contains the Euler--Lagrange terms
coming from varying curvature.  That is not a defect.  The defect of
\(f(K)g\) was precisely that the fourth-order derivative structure was hidden
and then lost in the homogeneous reduction.

\begin{proposition}[scale derivative]
Let \(M\) have dimension \(n=3\).  Under the constant conformal scaling
\(g_s=e^{2s}g\),
\[
        K_s=e^{-4s}K,
        \qquad
        d\mu_{g_s}=e^{3s}d\mu_g .
\]
Hence
\[
        \frac{d}{ds}\bigg|_{s=0}\F_\Phi[g_s]
        =
        \int_M\left(3\Phi(K)-4K\Phi'(K)\right)d\mu_g .
\]
Consequently the negative \(L^2\)-gradient has the correct high-curvature
conformal sign whenever
\[
        4K\Phi'(K)-3\Phi(K)>0.
\]
\end{proposition}

\begin{proof}
Substitution gives
\[
        \F_\Phi[g_s]
        =
        \int_M e^{3s}\Phi(e^{-4s}K)d\mu_g .
\]
Differentiating at \(s=0\) gives the displayed identity.  If \(E_\Phi\) is the
\(L^2\)-gradient and \(v=2g\), then
\[
        \int_M \langle E_\Phi,2g\rangle\,d\mu
        =
        \int_M(3\Phi-4K\Phi')\,d\mu .
\]
Thus the negative gradient expands the metric in the high-curvature conformal
mode precisely when \(4K\Phi'-3\Phi\) is positive.
\end{proof}

\section{Power-law regularizers}

Take
\[
        \Phi(K)=K^q .
\]
Then
\[
        4K\Phi'(K)-3\Phi(K)
        =
        (4q-3)K^q .
\]
Thus the conformal sign is correct exactly when
\[
        q>\frac34 .
\]

The same exponent is the scale threshold.  Under a length scale
\(L\),
\[
        |\Rm|\sim L^{-2},
        \qquad
        K=|\Rm|^2\sim L^{-4},
        \qquad
        d\mu\sim L^3,
\]
and therefore
\[
        \int K^q\,d\mu
        \sim
        L^{3-4q}.
\]
The functional is scale-critical at \(q=3/4\) and blows up on shrinking
curvature concentrations when \(q>3/4\).

\begin{remark}[correction of the exponent]
The exponent \(q>3/2\) is stronger than necessary for the scale and conformal
tests in dimension three.  The critical exponent for
\(\int |\Rm|^{2q}d\mu=\int K^q d\mu\) is \(q=3/4\).  The condition
\(q>3/2\) may be retained as a deliberately aggressive, non-minimal choice,
but it should not be called the minimal scale threshold.
\end{remark}

In a round high-curvature mode the contribution to \(K\) has the leading
form
\[
        \dt K\big|_{\mathrm{var}}
        =
        -\frac{\varepsilon}{3}(4q-3)K^{q+1}
        +\text{derivative and lower-order terms}.
\]
The Ricci-flow reaction has order \(K^{3/2}\).  Thus every
\[
        q>\frac12
\]
beats the cubic Ricci reaction by degree, and every \(q>3/4\) both beats it
and passes the scale/conformal test.

The power-law candidate is therefore
\[
        \boxed{
        \dt g
        =
        -2\Ric
        -
        \varepsilon\,\grad_{L^2}\int_{S^3} (|\Rm|^2+\delta)^q\,d\mu_g,
        \qquad
        q>\frac34,\quad \delta>0.
        }
\]
The parameter \(\delta>0\) makes the integrand smooth at \(K=0\) when \(q\)
is not an integer.

\section{Planck barrier}

The power-law energy punishes concentration, but it does not by itself encode
a finite ceiling.  A genuine curvature wall is obtained by restricting to
\[
        \U_P=\{g:\sup_M K<\KP\},
        \qquad
        \KP=\lP^{-4},
\]
and using
\[
        \B_{p,q}[g]
        =
        \int_M
        \frac{(K+\delta)^q}{(\KP-K)^p}\dmu,
        \qquad
        p>0,\quad q>\frac34,\quad \delta>0.
\]
The proposed Planck-barrier flow is
\[
        \boxed{
        \dt g
        =
        -2\Ric
        -
        \varepsilon\,\grad_{L^2}\B_{p,q}[g].
        }
\]

\begin{proposition}[barrier sign]
For
\[
        \Phi(K)=\frac{(K+\delta)^q}{(\KP-K)^p},
        \qquad 0\le K<\KP,
\]
one has
\[
        4K\Phi'(K)-3\Phi(K)
        =
        \Phi(K)
        \left(
             \frac{4qK}{K+\delta}
             +
             \frac{4pK}{\KP-K}
             -
             3
        \right).
\]
In particular, for large \(K\) and especially as \(K\nearrow\KP\), the
negative gradient gives an inward high-curvature drift, with leading blow-up
\[
        4K\Phi'(K)-3\Phi(K)
        \sim
        \frac{4pK(K+\delta)^q}{(\KP-K)^{p+1}}.
\]
Consequently the conformal contribution to \(K\) has leading sign
\[
        \dt K\big|_{\mathrm{bar}}
        \sim
        -
        \frac{4\varepsilon p}{3}
        \frac{K^2(K+\delta)^q}{(\KP-K)^{p+1}}
        \quad
        \text{as } K\nearrow\KP .
\]
\end{proposition}

\begin{proof}
Differentiate logarithmically:
\[
        \frac{\Phi'(K)}{\Phi(K)}
        =
        \frac{q}{K+\delta}
        +
        \frac{p}{\KP-K}.
\]
Multiplication by \(4K\) and subtraction of \(3\) give the formula.  The final
asymptotic follows from the pole at \(K=\KP\).  The factor \(K/3\) in the
drift of \(K\) is the conformal scaling relation in dimension three.
\end{proof}

\begin{remark}[choice of \(p\)]
For the pointwise wall, \(p>0\) already creates a singular barrier.  Taking
\(p>1\) or \(p>3/2\) is a stronger modelling choice.  It may be useful when one
wants the variational force to become extremely stiff before \(K\) reaches
\(\KP\), but it is not required by the elementary scale test.
\end{remark}

\section{Berger verification}

For the Berger family,
\[
        K=\frac{4(16-24u+11u^2)}{x^2},
        \qquad x=a^2,\quad u=\frac{c^2}{a^2}.
\]
The round point is \(u=1\), where
\[
        K=\frac{12}{x^2}.
\]
The unforced Berger calculation gives
\[
        \left.\frac{\dt K}{K^{3/2}}\right|_{u=1}
        =
        \frac{4}{\sqrt3}.
\]
Thus the Ricci reaction is positive cubic in \(|\Rm|\), equivalently
positive \(K^{3/2}\).

For the power-law variational forcing, the leading conformal high-curvature
term is
\[
        -\frac{\varepsilon}{3}(4q-3)K^{q+1}.
\]
For \(q>3/4\),
\[
        K^{q+1}\gg K^{3/2}
        \qquad (K\to\infty),
\]
so the proposed forcing passes the Berger round-collapse degree test.

For the Planck barrier,
\[
        \dt K\big|_{\mathrm{bar}}
        \sim
        -
        \frac{4\varepsilon p}{3}
        \frac{K^2(K+\delta)^q}{(\KP-K)^{p+1}},
        \qquad K\nearrow\KP,
\]
so the leading formal drift points away from the forbidden boundary
\(K=\KP\).

This calculation also explains why the Hopf fibre is not the missing
regularizing channel.  At \(c\to0\) with \(a\) fixed,
\[
        K\to \frac{64}{a^4}<\infty.
\]
The singular variable is \(a\), not \(c\).  The invariant \(K\) sees this
without imposing a preferred fibre direction.

\section{Why this is not another scalar multiplier}

It is tempting to replace the variational term by
\[
        h=2\lambda K^\alpha g.
\]
This is not the same mechanism.  It preserves only the zeroth-order conformal
part and hides the spatial derivative terms.  Once \(K\) is inhomogeneous,
the evolution of curvature includes Laplacians and Hessians of \(K^\alpha\).
Those terms are of the same scale as the zeroth-order damping at a neck.  The
homogeneous success is therefore not stable.

The variational construction is different.  The tensor
\[
        E_\Phi=\grad_{L^2}\int \Phi(|\Rm|^2)d\mu
\]
contains the derivative terms forced by the variation of curvature.  After
DeTurck gauge fixing, this belongs to the same general class as higher-order
curvature flows.  The model therefore pays the price necessary for PDE
honesty: fourth-order terms are present explicitly, not suppressed.

\section{Relation to second-order uniqueness}

One should not expect a universal second-order natural tensor to do this job.
Second-order divergence-free natural symmetric tensors are severely restricted
by Lovelock-type classification results.  In low dimensions this leaves only
metric and Einstein-type structures, not a new high-curvature damping channel.
This supports, but does not by itself prove, the conclusion that a genuine
local regularizer must leave the second-order class.

The present proposal leaves it by using a curvature-energy gradient.

\section{What is proved algebraically}

The construction gives an algebraically verified high-curvature regulator.
It does not, by itself, give global regularity or noncollapse.

\begin{theorem}[algebraic high-curvature verification]\label{thm:algebraic-verification}
Let \(M^3\) be compact and let \(K=|\Rm|^2\).  For the functional
\[
        \F_q[g]=\int_M (K+\delta)^q\,d\mu_g
\]
with \(q>3/4\), the negative \(L^2(d\mu)\)-gradient has the correct conformal
high-curvature sign and produces a leading round-mode drift of order
\(-K^{q+1}\), which dominates the Ricci reaction \(+K^{3/2}\).  For the
barrier functional
\[
        \B_{p,q}[g]
        =
        \int_M
        \frac{(K+\delta)^q}{(\KP-K)^p}\,d\mu_g
\]
with \(p>0\), the same sign becomes singular as \(K\nearrow\KP\).
\end{theorem}

\begin{proof}
The first assertion is the scale derivative computation with
\(\Phi(K)=(K+\delta)^q\).  For \(K\gg\delta\),
\[
        4K\Phi'-3\Phi
        =
        (4q-3)K^q+o(K^q),
\]
which is positive for \(q>3/4\).  The induced conformal drift in \(K\) has an
additional factor \(-\varepsilon K/3\), giving order \(-K^{q+1}\).  Since
\(q+1>3/2\), this dominates \(K^{3/2}\).  The barrier assertion follows from
\[
        4K\Phi'(K)-3\Phi(K)
        =
        \Phi(K)
        \left(
        \frac{4qK}{K+\delta}
        +
        \frac{4pK}{\KP-K}
        -3
        \right),
\]
whose second term diverges to \(+\infty\) as \(K\nearrow\KP\).
\end{proof}

\section{Why noncollapse is a different problem}

The curvature wall prevents a high-curvature channel.  It cannot, by itself,
prevent bounded-curvature collapse.  The same Hopf--Berger calculation already
shows the obstruction.

\begin{lemma}[bounded-curvature collapse is invisible to \(K\)-walls]\label{lem:hopf-collapse-obstruction}
Let
\[
        g_{1,c}=\sigma_1^2+\sigma_2^2+c^2\sigma_3^2
\]
be the Berger family on \(S^3\).  As \(c\downarrow0\),
\[
        \Vol(g_{1,c})\to0,\qquad \inj(g_{1,c})\to0,
\]
but
\[
        \sup_{S^3}|\Rm(g_{1,c})|^2
\]
remains bounded.  Hence no barrier depending only on the condition
\(K\nearrow\KP\) can be a complete noncollapse mechanism.
\end{lemma}

\begin{proof}
This is the \(u=z/x\to0\), \(x=1\) case of the curvature dictionary
\[
        |\Rm|^2
        =
        \frac{4(16-24u+11u^2)}{x^2}.
\]
With \(x=1\), the right-hand side remains bounded and tends to \(64\).
But the Hopf fibre length is proportional to \(c\), hence the injectivity
radius and volume collapse to zero.  Thus bounded-curvature collapse is a
separate channel.
\end{proof}

\begin{remark}[geometric and spectral meaning of the Hopf collapse]\label{rem:hopf-quotient-spectrum}
The limit \(c\downarrow0\) is a Cheeger--Gromov collapse along the free Hopf
\(S^1\)-action.  The limiting metric tensor has vertical kernel and descends
to the \(S^2\) base of the Hopf fibration; it is not a high-curvature limit.
This also explains why the first Laplace eigenvalue alone is not a reliable
collapse detector.  Fibre-invariant test functions have Rayleigh quotients
controlled by the base geometry, while vertical modes become expensive as
\(c^{-2}\).  Thus a useful spectral detector must be normalized or replaced by
a heat-kernel, Sobolev, systolic, or injectivity-radius quantity.
\end{remark}


\begin{proposition}[a heat-kernel detector sees fixed-base Hopf collapse]\label{prop:heat-kernel-hopf}
For \(r>0\), define
\[
        \mathcal H_r(g)
        =
        \sup_{x\in M,\;0<t\le r^2} t^{3/2}H_g(t,x,x),
\]
where \(H_g\) is the scalar heat kernel.  Along the fixed-base Hopf collapse
\[
        g_{a,c}=a^2(\sigma_1^2+\sigma_2^2)+c^2\sigma_3^2,
        \qquad a>0,\qquad c\downarrow0,
\]
one has
\[
        \mathcal H_r(g_{a,c})\to+\infty .
\]
More precisely,
\[
        \mathcal H_r(g_{a,c})
        \ge
        \frac{r^3}{\Vol(g_{a,c})}
        =
        \frac{r^3}{2\pi^2a^2c}.
\]
\end{proposition}

\begin{proof}
For any compact Riemannian manifold, the zero eigenfunction gives the
pointwise lower bound
\[
        H_g(t,x,x)\ge \frac{1}{\Vol(g)}
\]
for every \(t>0\).  Taking \(t=r^2\) in the definition of \(\mathcal H_r\)
therefore gives
\[
        \mathcal H_r(g)\ge \frac{r^3}{\Vol(g)}.
\]
For the Berger family used here,
\[
        \Vol(g_{a,c})=2\pi^2a^2c.
\]
Thus
\[
        \mathcal H_r(g_{a,c})
        \ge
        \frac{r^3}{2\pi^2a^2c}\to+\infty
        \qquad(c\downarrow0).
\]
\end{proof}

\begin{remark}[scope of the heat-kernel observation]\label{rem:heat-kernel-scope}
Proposition~\ref{prop:heat-kernel-hopf} is deliberately modest.  It proves
that a heat-kernel scale detector sees the fixed-base Hopf collapse, whereas
\(\lambda_1(-\Delta)\) alone does not.  It does not prove NC2 monotonicity for
\(\mathcal H_r\), and it does not replace the sharper adiabatic asymptotic
\[
        H_{g_{a,c}}(t,x,x)
        \sim
        \frac{1}{4\pi t}\frac{1}{2\pi c}
        \qquad(c^2\ll t\ll a^2),
\]
whose rigorous proof requires either the exact Berger spectrum or a quoted
adiabatic heat-kernel theorem.  Thus the heat kernel is a natural
non-additive collapse detector, not yet the missing noncollapse entropy.
\end{remark}


\begin{proposition}[Ricci lower bound from a wall gap]\label{prop:ricci-lower-from-wall-gap}
Let \(M^3\) be a Riemannian three-manifold.  If
\[
        K=|\Rm|^2\le K_0,
\]
then
\[
        \Ric
        \ge
        -\left(\frac12+\frac1{\sqrt3}\right)\sqrt{K_0}\,g .
\]
In particular, a wall gap
\[
        K\le \KP-\eta
\]
implies
\[
        \Ric
        \ge
        -\kappa_0 g,
        \qquad
        \kappa_0
        =
        \left(\frac12+\frac1{\sqrt3}\right)\sqrt{\KP-\eta}.
\]
\end{proposition}

\begin{proof}
In dimension three,
\[
        K=4|\Ric_0|^2+\frac13R^2.
\]
Thus
\[
        |\Ric_0|\le\frac{\sqrt K}{2},
        \qquad
        |R|\le \sqrt{3K}.
\]
Every eigenvalue of the traceless symmetric tensor \(\Ric_0\) is bounded below
by \(-|\Ric_0|\), while
\[
        \Ric=\Ric_0+\frac R3g.
\]
Therefore
\[
        \lambda_{\min}(\Ric)
        \ge
        -|\Ric_0|-\frac{|R|}{3}
        \ge
        -\left(\frac12+\frac1{\sqrt3}\right)\sqrt K.
\]
The stated wall-gap estimate follows by inserting \(K\le \KP-\eta\).
\end{proof}

\begin{corollary}[heat-kernel coercivity under a wall gap]\label{cor:heat-kernel-coercivity}
Assume on a three-manifold that
\[
        K\le\KP-\eta .
\]
Then the constants in the local on-diagonal heat-kernel comparison estimates
may be taken to depend only on \(\KP,\eta\) and the chosen scale.  In
particular, for \(0<\rho\le r\), standard Li--Yau/Davies type comparison under
the Ricci lower bound of Proposition~\ref{prop:ricci-lower-from-wall-gap}
gives a constant \(c=c(\KP,\eta,r)>0\) such that
\[
        H_g(\rho^2,x,x)
        \ge
        \frac{c}{\Vol B(x,\rho)}.
\]
Consequently, if
\[
        \mathcal H_r(g)\le A,
\]
then
\[
        \Vol B(x,\rho)
        \ge
        c\,A^{-1}\rho^3
        \qquad
        (0<\rho\le r).
\]
\end{corollary}

\begin{proof}
By Proposition~\ref{prop:ricci-lower-from-wall-gap}, the wall gap gives a
uniform lower Ricci bound
\[
        \Ric\ge-\kappa_0g,
        \qquad
        \kappa_0
        =
        \left(\frac12+\frac1{\sqrt3}\right)\sqrt{\KP-\eta}.
\]
The local Li--Yau/Davies heat-kernel comparison theory under this lower Ricci
bound gives the displayed on-diagonal lower estimate, with constants depending
only on dimension, \(\kappa_0\), and the scale \(r\).  On the other hand,
\(\mathcal H_r(g)\le A\) gives
\[
        H_g(\rho^2,x,x)\le A\rho^{-3}.
\]
Combining the two inequalities yields
\[
        \frac{c}{\Vol B(x,\rho)}
        \le
        A\rho^{-3},
\]
hence
\[
        \Vol B(x,\rho)\ge c\,A^{-1}\rho^3 .
\]
\end{proof}

\begin{remark}[meaning of the heat-kernel direction]\label{rem:heat-kernel-direction}
Corollary~\ref{cor:heat-kernel-coercivity} supplies an NC3-type coercivity
direction for the heat-kernel scale detector, provided a wall gap is already
known.  It does not prove NC2 propagation of \(\mathcal H_r\) along the
curvature-wall flow.  The remaining noncollapse problem is therefore the
propagation estimate, or a replacement for it:
\[
        \mathcal H_r(g(t))
        \le
        C(T,\KP,\eta,r)\,\mathcal H_r(g(0)).
\]
The known spectra of Berger spheres make this propagation question directly
testable in the homogeneous sector, but such spectral probes are numerical
evidence unless the normalization and the running supremum in \(\mathcal H_r\)
are treated as part of a separate argument.
\end{remark}


Therefore a complete replacement for surgery cannot be only a local
high-curvature wall.  It must also contain a noncollapse principle.

\section{NC1--NC3: the required noncollapse package}

The noncollapse closure is encoded by three properties.  These are stated here
explicitly because they are not consequences of the local curvature-wall
calculation.

\begin{definition}[NC1--NC3]\label{def:NC}
Let
\[
        \Phi(K)=\frac{(K+\delta)^q}{(\KP-K)^p},
        \qquad q>3/4,\quad p>0,\quad \delta>0,
\]
on
\[
        \U_P=\{g:\sup_MK<\KP\}.
\]
A Perelman-type noncollapse functional for the barrier problem is a family
\(\Ncal(g,f,\tau)\), with \(f\in C^\infty(M)\), \(\tau>0\), satisfying:

\begin{enumerate}[label=\textbf{NC\arabic*.}]
\item \emph{Correct scale weight.}  Under
\(g\mapsto\theta^2g\), \(\tau\mapsto\theta^2\tau\), the barrier-coupling term
has the unique power \(\tau^{2q}\) in the pure model \(\Phi(K)=K^q\).

\item \emph{Monotonicity.}  Along the relevant coupled flow of
\((g,f,\tau)\), the minimized entropy is monotone in the Perelman direction.

\item \emph{Coercivity implies noncollapse.}  If the entropy is bounded below
and \(K\le r^{-4}\) on \(B(x,r)\), with \(r^{-4}<\KP\), then
\[
        \Vol B(x,r)\ge \kappa(A,r)\,r^3 .
\]
The constant is allowed to depend on the fixed scale \(r\), because \(\KP\)
breaks exact scale invariance.
\end{enumerate}
\end{definition}

The first scale-correct candidate is
\[
        \Wcal_\Phi(g,f,\tau)
        =
        \int_M
        \left[
        \tau(R+|\nabla f|^2)
        +
        \varepsilon\tau^{2q}\Phi(K)
        +
        f-3
        \right]
        (4\pi\tau)^{-3/2}e^{-f}\,d\mu .
\]
It is important to call this only a \emph{candidate}.  The power
\(\tau^{2q}\) is forced by scaling, but scaling does not imply monotonicity.
For \(\Phi=K^q\), NC1 is exact:
\[
        (\theta^2\tau)^{2q}\Phi(\theta^{-4}K)
        =
        \theta^{4q}\tau^{2q}\theta^{-4q}K^q
        =
        \tau^{2q}K^q .
\]
NC3 is soft once NC2 is known: the usual collapsed-ball test function gives
\(\Wcal(g,f^\ast,\tau)\le\log\kappa_0+C\), and the extra term is bounded on
\(K\le r^{-4}<\KP\).  Thus \(\mu_\Phi\to-\infty\) as \(\kappa_0\to0\),
contradicting an entropy lower bound.  The entire hard content is NC2.  The
Berger calculation below shows that the natural eigenvalue version of this
candidate does not satisfy NC2.

\section{The NC2 obstruction for the unweighted barrier flow}

Write
\[
        \mathcal P_\Phi[g]=\int_M\Phi(K)\,d\mu_g,
        \qquad K=|\Rm|^2,
\]
and define \(E_\Phi\) by
\[
        D\mathcal P_\Phi[g](h)
        =
        \int_M \langle E_\Phi,h\rangle\,d\mu .
\]
With the curvature-variation conventions used here, the formal first-variation
tensor has the Streets-type structure
\[
        (E_\Phi)_{ij}
        =
        \frac12\Phi g_{ij}
        -
        2\Phi'(K)R_i{}^{abc}R_{jabc}
        +
        4\nabla^a\nabla^b\bigl(\Phi'(K)R_{iabj}\bigr).
\]
For \(\Phi(K)=K\) this is the usual \(L^2\)-curvature-gradient structure,
up to the sign and index conventions of the chosen curvature tensor.  In this
paper the theorem-level consequences that use this tensor are checked on the
homogeneous Berger family by symmetric criticality and by direct differentiation
of the reduced functional.  The full tensorial expression is recorded to explain
the analytic structure of the flow; a complete PDE theorem would require an
independent convention-by-convention comparison with the standard
\(L^2\)-curvature-flow formula.

The bare barrier flow is
\[
        \dt g=-2\Ric-\varepsilon E_\Phi .
\]

Perelman's entropy mechanism, however, is not naturally expressed in the
unweighted inner product.  It uses the weighted measure
\[
        \rho\,d\mu=(4\pi\tau)^{-3/2}e^{-f}\,d\mu .
\]
Define the \(f\)-weighted Euler tensor \(E_\Phi^f\) by
\[
        D\int_M\Phi(K)\rho\,d\mu
        =
        \int_M\langle E_\Phi^f,h\rangle\,\rho\,d\mu
\]
with \(f\) and \(\tau\) fixed during the metric variation.  With the same convention, the corresponding weighted expression is
\[
        (E_\Phi^f)_{ij}
        =
        \frac12\Phi g_{ij}
        -
        2\Phi'(K)R_i{}^{abc}R_{jabc}
        +
        4e^f\nabla^a\nabla^b
        \left(e^{-f}\Phi'(K)R_{iabj}\right).
\]
Let
\[
        Q_{iabj}=\Phi'(K)R_{iabj}.
\]
The discrepancy between the two expressions is the weighted-divergence
correction
\[
        (E_\Phi^f-E_\Phi)_{ij}
        =
        4\left[
        e^f\nabla^a\nabla^b(e^{-f}Q_{iabj})
        -
        \nabla^a\nabla^bQ_{iabj}
        \right].
\]
Expanding the weighted derivatives gives
\[
\begin{aligned}
        (E_\Phi^f-E_\Phi)_{ij}
        =
        4\big[
        &-(\nabla^a f)\nabla^bQ_{iabj}
        -(\nabla^b f)\nabla^aQ_{iabj}  \\
        &+
        \big((\nabla^a f)(\nabla^b f)-\nabla^a\nabla^b f\big)Q_{iabj}
        \big].
\end{aligned}
\]

\begin{lemma}[no pointwise sign for the discrepancy]\label{lem:no-sign}
The tensor \(E_\Phi^f-E_\Phi\) has no intrinsic pointwise sign.
\end{lemma}

\begin{proof}
At a point where \(Q\ne0\), choose normal coordinates and choose \(f\) so that
\(\nabla f=0\) at that point while \(\nabla^2f\) is arbitrary.  Then
\[
        (E_\Phi^f-E_\Phi)_{ij}
        =
        -4(\nabla^a\nabla^b f)Q_{iabj}.
\]
Contracting with a symmetric tensor \(S^{ij}\) gives
\[
        S^{ij}(E_\Phi^f-E_\Phi)_{ij}
        =
        -4S^{ij}(\nabla^a\nabla^b f)Q_{iabj}.
\]
The Hessian of \(f\) may be chosen with either sign in any nonzero contraction
direction.  Thus the expression can be positive or negative.
\end{proof}

Consequently the original unweighted barrier flow has no automatic
Perelman-type NC2 theorem.  One must either prove a separate absorption
estimate for this discrepancy or change the flow.


\section{Berger falsification of the naive eigenvalue candidate}

The preceding obstruction \(E_\Phi^f-E_\Phi\) is not the only problem.  The
most natural eigenvalue candidate fails even in the homogeneous sector where
that obstruction disappears.

Define
\[
        \lambda_\Phi(g)
        =
        \lambda_1\!\left(-4\Delta_g+R+\varepsilon\Phi(K)\right).
\]
On a homogeneous Berger metric, \(R\) and \(K\) are spatially constant.  Hence
the ground state of \(-4\Delta_g+R+\varepsilon\Phi(K)\) is the constant
function and
\[
        \lambda_\Phi(g)=R+\varepsilon\Phi(K)
\]
exactly.  Thus \(f\) is constant, \(\nabla f=0\), and the weighted-gradient
ambiguity is absent.  This is the most favourable sector for a Perelman-type
NC2 candidate.

Let
\[
        g_{a,c}=a^2(\sigma_1^2+\sigma_2^2)+c^2\sigma_3^2,\qquad
        x=a^2,\quad z=c^2,\quad u=z/x .
\]
The Berger curvature dictionary is
\[
        K=\frac{4(16x^2-24xz+11z^2)}{x^4},\qquad
        R=\frac8x-\frac{2z}{x^2},
\]
and the unforced Ricci-flow ODE is
\[
        \dot x_0=-8+\frac{4z}{x},\qquad
        \dot z_0=-\frac{4z^2}{x^2}.
\]
At \(u=1\), \(R=6/x\), \(K=12/x^2\), and
\[
        C_3^{\mathrm B}
        =
        \frac{\dot K}{K^{3/2}}
        =
        \frac4{\sqrt3}.
\]


\begin{proposition}[explicit Berger form of \(E_\Phi\)]\label{prop:Ephi}
For any smooth \(\Phi\), on the Berger family
\[
        g_{a,c}=a^2(\sigma_1^2+\sigma_2^2)+c^2\sigma_3^2,
        \qquad x=a^2,\quad z=c^2,
\]
with
\[
        K=\frac{4(16x^2-24xz+11z^2)}{x^4},
\]
the symmetric-criticality reduction of \(E_\Phi=\grad_{L^2}\int\Phi(K)d\mu\)
is
\[
\boxed{
        (E_\Phi)_{11}
        =
        -\frac{8\Phi'(K)(8x^2-18xz+11z^2)}{x^3}
        +
        \frac{x}{2}\Phi(K),
}
\]
and
\[
\boxed{
        (E_\Phi)_{33}
        =
        -\frac{8z^2\Phi'(K)(12x-11z)}{x^4}
        +
        \frac{z}{2}\Phi(K).
}
\]
At the round point \(z=x\), both components reduce to
\[
        (E_\Phi)_{11}=(E_\Phi)_{33}
        =
        -\frac{8\Phi'(K)}{x}
        +
        \frac{x}{2}\Phi(K),
\]
as required by the enhanced isotropy of the round metric.
\end{proposition}

\begin{proof}
The reduced functional is
\[
        \mathcal P_\Phi(x,z)=\Phi(K(x,z))\Vol(x,z).
\]
By symmetric criticality,
\[
        (E_\Phi)_{11}
        =
        \frac{x^2}{2\Vol}\partial_x\mathcal P_\Phi,
        \qquad
        (E_\Phi)_{33}
        =
        \frac{z^2}{\Vol}\partial_z\mathcal P_\Phi .
\]
Since
\[
        K_x=
        -\frac{16(8x^2-18xz+11z^2)}{x^5},
        \qquad
        K_z=
        -\frac{8(12x-11z)}{x^4},
\]
and
\[
        \frac{\partial_x\Vol}{\Vol}=\frac1x,\qquad
        \frac{\partial_z\Vol}{\Vol}=\frac1{2z},
\]
the displayed formulas follow immediately.  For \(\Phi(K)=K\) they reduce to
the components used below.
\end{proof}

\begin{corollary}[Planck-wall components and their signs]\label{cor:planck-component-signs}
For the physical wall
\[
        \Phi(K)=\frac{(K+\delta)^q}{(\KP-K)^p},
        \qquad p>0,\quad q>0,\quad \delta>0,
\]
put
\[
        A(K)=\frac{\Phi'(K)}{\Phi(K)}
        =
        \frac{q}{K+\delta}
        +
        \frac{p}{\KP-K}.
\]
Then Proposition~\ref{prop:Ephi} becomes
\[
\boxed{
        (E_\Phi)_{11}
        =
        \Phi(K)
        \left[
        \frac{x}{2}
        -
        \frac{8}{x}(11u^2-18u+8)A(K)
        \right],
}
\]
and
\[
\boxed{
        (E_\Phi)_{33}
        =
        \Phi(K)
        \left[
        \frac{ux}{2}
        -
        \frac{8u^2}{x}(12-11u)A(K)
        \right].
}
\]
Since
\[
        11u^2-18u+8>0
\]
for all real \(u\), the horizontal component is negative whenever
\[
        x^2<16(11u^2-18u+8)A(K).
\]
In particular, as \(K\nearrow\KP\), one has \(A(K)\to+\infty\), and therefore
\[
        (E_\Phi)_{11}<0
\]
near the wall.

For the vertical component, if \(u>12/11\), then \(12-11u<0\), and hence
\[
        (E_\Phi)_{33}>0.
\]
If \(0<u<12/11\), then
\[
        (E_\Phi)_{33}<0
        \quad\Longleftrightarrow\quad
        x^2<16u(12-11u)A(K).
\]
Consequently, near the wall,
\[
        (E_\Phi)_{33}<0\quad (0<u<12/11),
        \qquad
        (E_\Phi)_{33}>0\quad (u>12/11).
\]
Thus the wall has a definite leading sign structure on the Berger sectors:
it expands horizontal directions near the wall, while the vertical reaction
changes sign at \(u=12/11\).
\end{corollary}


\begin{proposition}[reduced Berger dynamics for a general wall]\label{prop:berger-general-wall-ode}
On the Berger family, the curvature-wall equation
\[
        \dot g=-2\Ric-\varepsilon E_\Phi
\]
reduces to
\[
\boxed{
        \dot x
        =
        -8+4u
        +
        \frac{8\varepsilon\Phi'(K)}{x}(8-18u+11u^2)
        -
        \frac{\varepsilon x}{2}\Phi(K),
}
\]
and
\[
\boxed{
        \dot z
        =
        -4u^2
        +
        \frac{8\varepsilon u^2\Phi'(K)}{x}(12-11u)
        -
        \frac{\varepsilon ux}{2}\Phi(K).
}
\]
Consequently, for \(u=z/x\),
\[
\boxed{
        \dot u
        =
        \frac{8u(1-u)}{x}
        \left(
        1+\frac{2\varepsilon\Phi'(K)}{x}(11u-4)
        \right).
}
\]
In the anisotropy equation the zeroth-order \(\Phi(K)\)-terms cancel; only the
slope \(\Phi'(K)\) of the wall affects the fibre ratio.
\end{proposition}

\begin{proof}
The \(x,z\)-system is obtained by substituting Proposition~\ref{prop:Ephi} into
\[
        \dot x=\dot x_0-\varepsilon(E_\Phi)_{11},\qquad
        \dot z=\dot z_0-\varepsilon(E_\Phi)_{33}.
\]
Since \(z=ux\), one has
\[
        \dot u=\frac{\dot z-u\dot x}{x}.
\]
Substitution gives the displayed formula.  The cancellation of the
\(\Phi(K)\)-terms follows from their pure scaling character:
\[
        -\frac{\varepsilon ux}{2}\Phi(K)
        -
        u\left(-\frac{\varepsilon x}{2}\Phi(K)\right)=0 .
\]
\end{proof}

\begin{corollary}[Berger cone and Hopf-sector warning]\label{cor:berger-cone}
Assume \(\varepsilon>0\) and \(\Phi'(K)>0\).  Then, on the Berger sector,
\[
        \frac4{11}\le u\le M,\qquad M>1,
\]
is forward invariant.  More precisely, \(u=1\) is invariant; for
\(\frac4{11}\le u<1\) one has \(\dot u>0\), and for \(u>1\) one has
\(\dot u<0\).  In the lower Hopf sector \(0<u<4/11\), the wall can instead
push \(u\) downward whenever
\[
\boxed{
        \frac{2\varepsilon\Phi'(K)}{x}(4-11u)>1 .
}
\]
Thus a \(K\)-wall controls high curvature but need not oppose every
bounded-curvature fibre-collapse direction.
\end{corollary}

\begin{proof}
The sign of \(\dot u\) is the sign of
\[
        u(1-u)
        \left(
        1+\frac{2\varepsilon\Phi'(K)}{x}(11u-4)
        \right).
\]
If \(u\ge4/11\), the parenthesis is at least \(1\).  The stated signs follow.
If \(0<u<4/11\), the parenthesis is negative precisely under the displayed
inequality.
\end{proof}

\begin{corollary}[curvature, scale, and volume consequences]\label{cor:berger-scale-volume}
Along the Berger family,
\[
        K=\frac{4(16-24u+11u^2)}{x^2}.
\]
Hence a lower bound \(x\ge x_0>0\) gives a curvature bound only together with
an upper bound for \(u\).  In particular an \(a\)-barrier \(a^2=x\ge x_0\) is
not equivalent to the invariant curvature wall \(K<\KP\), since \(u\to\infty\)
at fixed \(x\) makes \(K\to\infty\).

The volume satisfies
\[
        \Vol(g_{a,c})=2\pi^2x\sqrt z
\]
and
\[
\boxed{
        \frac{d}{dt}\log\Vol
        =
        \frac{-8+2u}{x}
        +
        \frac{4\varepsilon\Phi'(K)}{x^2}(16-24u+11u^2)
        -
        \frac{3\varepsilon}{4}\Phi(K).
}
\]
On the round locus \(u=1\), this becomes
\[
\boxed{
        \frac{d}{dt}\log\Vol
        =
        -\frac6x
        +
        \frac{12\varepsilon\Phi'(K)}{x^2}
        -
        \frac{3\varepsilon}{4}\Phi(K).
}
\]
Volume normalization therefore does not by itself remove Hopf-fibre collapse:
at fixed volume \(x\sqrt z=\mathrm{const}\), the limit \(z\downarrow0\) is still
a fibre collapse, with \(x\to\infty\).
\end{corollary}

\begin{proof}
The curvature formula is the Berger dictionary.  The statement about \(x\) and
\(u\) is immediate from the displayed expression for \(K\).  For the volume,
\[
        \frac{d}{dt}\log\Vol
        =
        \frac{\dot x}{x}
        +
        \frac12\frac{\dot z}{z}.
\]
Substituting Proposition~\ref{prop:berger-general-wall-ode} gives the displayed
formula.
\end{proof}

\begin{corollary}[normalization, scalar curvature, and traceless anisotropy]\label{cor:berger-scalar-anisotropy}
Volume normalization changes only the common scale in the Berger reduction.
Indeed, replacing a velocity \(v\) by
\[
        v^{\mathrm{norm}}
        =
        v-\frac13\overline{\tr_g v}\,g
\]
changes
\[
        \dot x\mapsto \dot x-\rho x,\qquad
        \dot z\mapsto \dot z-\rho z
\]
for a scalar \(\rho(t)\), and therefore leaves
\[
        u=\frac{z}{x}
\]
unchanged:
\[
\boxed{
        \dot u_{\mathrm{norm}}=\dot u_{\mathrm{unnorm}}.
}
\]
Thus volume normalization cannot by itself remove Hopf-fibre collapse; it
changes the common scale but not the anisotropy equation.

The scalar curvature and traceless Ricci size are
\[
        R=\frac{8-2u}{x},
        \qquad
\boxed{
        |\Ric_0|^2=\frac{32(1-u)^2}{3x^2}.
}
\]
Along the general Berger wall-flow,
\[
\boxed{
        \dot R
        =
        \frac{
        \varepsilon\Phi(K)x^2(4-u)
        +
        16\varepsilon\Phi'(K)(3u-4)(11u^2-16u+8)
        +
        8x(3u^2-8u+8)
        }{x^3}.
}
\]
At \(u=1\), this reduces to
\[
\boxed{
        \dot R
        =
        \frac{24}{x^2}
        +
        \frac{3\varepsilon\Phi(K)}{x}
        -
        \frac{48\varepsilon\Phi'(K)}{x^3}.
}
\]
The dimensionless shape anisotropy decreases in the protected cone:
if \(u\ge4/11\), \(u\ne1\), \(\varepsilon>0\), and \(\Phi'(K)>0\), then
\[
\boxed{
        \frac{d}{dt}(1-u)^2<0.
}
\]
The absolute quantity \(|\Ric_0|^2\), however, can still grow if the scale
\(x\) decreases fast enough.  Thus the Berger dynamics separates shape
rounding from absolute curvature growth.
\end{corollary}

\begin{proof}
The normalization statement follows because subtracting a scalar multiple of
the metric subtracts the same logarithmic rate from \(x\) and \(z\).  The
formula for \(R\) is the Berger curvature dictionary, and the identity
\[
        K=4|\Ric_0|^2+\frac13R^2
\]
gives the displayed expression for \(|\Ric_0|^2\).  Differentiating
\(R=(8-2u)/x\) using Proposition~\ref{prop:berger-general-wall-ode} gives the
formula for \(\dot R\).  Finally,
\[
        \frac{d}{dt}(1-u)^2
        =
        -\frac{16u(1-u)^2}{x}
        \left(
        1+\frac{2\varepsilon\Phi'(K)}{x}(11u-4)
        \right),
\]
whose sign is negative in the stated cone.
\end{proof}



\begin{proposition}[exact Berger falsification for \(\Phi(K)=K\)]
For the pure quartic coupling \(\Phi(K)=K\), the reduced Euler tensor gives
\[
        (E_\Phi)_{11}
        =
        \frac{2(-16x^2+48xz-33z^2)}{x^3},
        \qquad
        (E_\Phi)_{33}
        =
        \frac{2z(16x^2-72xz+55z^2)}{x^4}.
\]
Along the forced Berger system
\[
        \dot x=\dot x_0-\varepsilon(E_\Phi)_{11},\qquad
        \dot z=\dot z_0-\varepsilon(E_\Phi)_{33},
\]
one has
\[
\boxed{
        \dot\lambda_\Phi(x,u,\varepsilon)
        =
        \frac{8c_0(u)}{x^2}
        -
        \frac{4c_1(u)}{x^3}\varepsilon
        -
        \frac{16c_2(u)}{x^4}\varepsilon^2,
}
\]
where
\[
        c_0(u)=3u^2-8u+8,
\]
\[
        c_1(u)=143u^3-436u^2+464u-192,
\]
\[
        c_2(u)=1331u^4-3696u^3+3648u^2-1536u+256 .
\]
At the round point \(u=1\),
\[
\boxed{
        \dot\lambda_\Phi\big|_{u=1}
        =
        \frac{12(2x-\varepsilon)(x+4\varepsilon)}{x^4}.
}
\]
\end{proposition}

\begin{proof}
The reduced functional is
\[
        \mathcal P_\Phi(x,z)=\Phi(K(x,z))\,\Vol(x,z).
\]
By symmetric criticality,
\[
        (E_\Phi)_{11}
        =
        \frac{x^2}{2\Vol}\,\partial_x\mathcal P_\Phi,
        \qquad
        (E_\Phi)_{33}
        =
        \frac{z^2}{\Vol}\,\partial_z\mathcal P_\Phi .
\]
For \(\Phi=K\), direct differentiation gives the displayed components of
\(E_\Phi\).  Substituting them into the forced ODE and differentiating
\[
        \lambda_\Phi=R+\varepsilon K
\]
gives the stated formula.  The round specialization \(z=x\) simplifies to
\[
        \dot\lambda_\Phi
        =
        \frac{24x^2+84x\varepsilon-48\varepsilon^2}{x^4}
        =
        \frac{12(2x-\varepsilon)(x+4\varepsilon)}{x^4}.
\]
\end{proof}

\begin{corollary}[the naive NC2 candidate is false]
For every fixed \(\varepsilon>0\),
\[
        \dot\lambda_\Phi\big|_{u=1}<0
        \qquad\text{whenever}\qquad
        x<\frac{\varepsilon}{2}.
\]
Thus \(\lambda_\Phi=R+\varepsilon K\) is not monotone along the barrier flow.
This is not a failure of an estimate; it is an exact sign reversal.
\end{corollary}

\begin{remark}[the sign reversal is not only round]
The polynomial \(c_0(u)=3u^2-8u+8\) is positive for all \(u>0\), since its
discriminant is negative.  The polynomial \(c_2(u)\) has no positive root by
Sturm's test and \(c_2(0)=256>0\).  Hence the negative
\(-16c_2(u)\varepsilon^2/x^4\) term eventually dominates for every fixed fibre
ratio \(u>0\).  The round calculation is only the simplest visible instance.
\end{remark}

\begin{proposition}[the Planck wall worsens the sign]
Let
\[
        \Phi(K)=\frac{(K+\delta)^q}{(\KP-K)^p},
        \qquad p>0,\quad q>0,\quad \delta>0 .
\]
At the round point, along the same unweighted barrier flow,
\[
        \dot\lambda_\Phi\longrightarrow -\infty
        \qquad\text{as}\qquad K\nearrow\KP
\]
for every fixed \(\varepsilon>0\).
\end{proposition}

\begin{proof}
At the round point,
\[
        R=\frac6x,\qquad K=\frac{12}{x^2}.
\]
For a pure scaling variation,
\[
        \dot x
        =
        -4
        +
        \frac{\varepsilon x}{6}
        \bigl(4K\Phi'(K)-3\Phi(K)\bigr).
\]
Since
\[
        \dot R=-\frac{R}{x}\dot x,\qquad
        \dot K=-\frac{2K}{x}\dot x,
\]
we get
\[
\begin{aligned}
        \dot\lambda_\Phi
        &=
        \dot R+\varepsilon\Phi'(K)\dot K                                      \\
        &=
        \frac{4R}{x}
        +
        \frac{8\varepsilon K\Phi'(K)}{x}
        -
        \varepsilon
        \left(
        \frac{R}{6}
        +
        \frac{\varepsilon K\Phi'(K)}{3}
        \right)
        \bigl(4K\Phi'(K)-3\Phi(K)\bigr).
\end{aligned}
\]
Write \(s=\KP-K\).  As \(s\downarrow0\),
\[
        \Phi(K)\sim C s^{-p},\qquad
        \Phi'(K)\sim pC s^{-p-1},
\]
with \(C=(\KP+\delta)^q>0\).  Therefore
\[
        4K\Phi'(K)-3\Phi(K)\sim 4pKC\,s^{-p-1},
\]
and the leading term in \(\dot\lambda_\Phi\) is
\[
        -\frac{\varepsilon^2K\Phi'(K)}{3}
        \bigl(4K\Phi'(K)-3\Phi(K)\bigr)
        \sim
        -\frac{4\varepsilon^2p^2K^2C^2}{3}\,s^{-2p-2}.
\]
All positive terms have lower order \(s^{-p-1}\).  Hence
\(\dot\lambda_\Phi\to-\infty\).
\end{proof}

The conclusion is sharp:
\[
        \boxed{\lambda_\Phi=R+\varepsilon\Phi(K)\ \text{is not an NC2 functional}.}
\]
The curvature wall still functions as a local metric damping term, but the
eigenvalue/entropy route \(R+\varepsilon\Phi(K)\) is closed.  Any future
noncollapse functional must be structurally different and must first pass this
Berger test.


\section{No-go for additive matching \texorpdfstring{\(K\)}{K}-potential repairs}

The preceding Berger falsification closes the most direct eigenvalue route.
A broader additive class also fails.  The following result is deliberately
stated only for the class that is actually tested: an additive \(K\)-potential
is inserted into the Perelman functional, and the reduced metric forcing is
generated by the same added scalar through symmetric criticality.  It does not
claim a no-go theorem for arbitrary smooth functions \(H(\tau,K)\).

On the round Berger locus \(g_x=xg_{\mathrm{round}}\),
\[
        R=\frac6x,\qquad K=\frac{12}{x^2}.
\]
Let
\[
        \Ncal_F
        =
        \Wcal_{\mathrm{Perelman}}
        +
        \sigma\varepsilon\tau^\alpha K,
        \qquad
        \sigma\in\{\pm1\},\quad \alpha\in\mathbb R .
\]
At the Berger level, the Ricci part is still the Ricci-flow part; only the
added \(K\)-term contributes its matching gradient forcing.  This avoids
confusing the tensor inside Perelman's monotonicity square,
\[
        \Ric+\nabla^2 f-\frac1{2\tau}g,
\]
with the actual Ricci-flow velocity \(-2\Ric\).

\begin{theorem}[Berger no-go for additive matching \(K\)-potentials]
For the pure representative \(\Phi(K)=K\), at the round point and with
\(\tau=1\),
\[
        \frac{d}{dt}\Ncal_F
        =
        \frac{3N(x)}{2x^4},
\]
where
\[
        N(x)
        =
        x^2(x-4)^2
        +
        \sigma\varepsilon(2-8\alpha)x^2
        +
        56\sigma\varepsilon x
        -
        32\varepsilon^2 .
\]
Consequently, for every \(\varepsilon>0\), every sign \(\sigma\), and every
exponent \(\alpha\), the derivative of \(\Ncal_F\) is negative for some smooth
round Berger initial data.  In particular,
\[
        N(0)=-32\varepsilon^2<0 .
\]
Thus no member of this additive matching \(K\)-potential class satisfies NC2.
\end{theorem}

\begin{proof}
The first term \(x^2(x-4)^2\) is the classical Perelman square in the round
reduction.  The remaining terms are obtained by differentiating the added
scalar \(\sigma\varepsilon\tau^\alpha K\) along the reduced matching-gradient
forcing and then setting \(z=x\), \(\tau=1\).

The displayed formula gives
\[
        N(0)=-32\varepsilon^2,
\]
which is independent of \(\sigma\) and \(\alpha\).  By continuity there exists
\(x_0>0\) such that \(N(x)<0\) for \(0<x<x_0\).  Choosing a smooth round
initial metric with such an \(x\) gives
\[
        \left.\frac{d}{dt}\Ncal_F\right|_{t=0}<0 .
\]
Monotonicity of an NC2 functional must hold for all smooth initial data in the
tested class.  Hence the additive matching \(K\)-potential class fails the
Berger test.
\end{proof}

\begin{theorem}[analytic near-wall no-go for Planck-barrier additive repairs]
Let
\[
        \Phi(K)=\frac{(K+\delta)^q}{(\KP-K)^p},
        \qquad
        p>0,\quad q>0,\quad \delta>0 .
\]
Consider the additive matching-potential class
\[
        \Ncal_{\Phi,\alpha}
        =
        \Wcal_{\mathrm{Perelman}}
        +
        \sigma\varepsilon\tau^\alpha \Phi(K),
        \qquad
        \sigma\in\{\pm1\},\quad \alpha\in\mathbb R,
\]
with the reduced metric forcing generated by the same added scalar.  At the
round Berger point, with \(\tau=1\), write
\[
        s=\KP-K,\qquad x_P=\sqrt{\frac{12}{\KP}} .
\]
Then, as \(s\downarrow0\),
\[
\boxed{
        \frac{d}{dt}\Ncal_{\Phi,\alpha}
        =
        -\frac{192\,p^2(\KP+\delta)^{2q}}{x_P^4}
        \,\varepsilon^2\,s^{-2p-2}
        +
        O_{\varepsilon,\sigma,\alpha}(s^{-p-1}) .
}
\]
In particular,
\[
        \frac{d}{dt}\Ncal_{\Phi,\alpha}\longrightarrow -\infty
        \qquad\text{as}\qquad K\nearrow\KP
\]
for every fixed \(\varepsilon\ne0\), independent of the sign \(\sigma\) and of
the exponent \(\alpha\).
\end{theorem}

\begin{proof}
On the round Berger locus, the classical part of Perelman's functional gives
the usual square.  Differentiating the additive potential along the
matching-gradient reduced forcing and organizing the result by powers of
\(\varepsilon\) gives
\[
        \frac{d}{dt}\Ncal_{\Phi,\alpha}
        =
        \frac{3(x-4)^2}{2x^2}
        +
        \sigma\varepsilon A_{\Phi,\alpha}(x)
        +
        \varepsilon^2
        \frac{12\Phi'(K)\bigl(\Phi(K)x^2-16\Phi'(K)\bigr)}{x^4},
\]
where \(A_{\Phi,\alpha}(x)=O(s^{-p-1})\) as \(s=\KP-K\downarrow0\).  The exact
form of \(A_{\Phi,\alpha}\) depends on the sign and \(\tau\)-weight, but it is
linear in the coupling and lower order at the wall.

Now
\[
        \Phi(K)\sim(\KP+\delta)^q s^{-p},
\]
and
\[
        \Phi'(K)
        =
        q(K+\delta)^{q-1}s^{-p}
        +
        p(K+\delta)^q s^{-p-1}
        \sim
        p(\KP+\delta)^q s^{-p-1}.
\]
Hence
\[
        \Phi(K)\Phi'(K)=O(s^{-2p-1}),
        \qquad
        \Phi'(K)^2
        \sim
        p^2(\KP+\delta)^{2q}s^{-2p-2}.
\]
Therefore the quadratic term has leading part
\[
        \varepsilon^2
        \frac{12\Phi'(K)(\Phi(K)x^2-16\Phi'(K))}{x^4}
        \sim
        -\frac{192p^2(\KP+\delta)^{2q}}{x_P^4}
        \varepsilon^2s^{-2p-2}.
\]
This dominates both the classical square and the linear coupling term.  The
leading coefficient is strictly negative and contains \(\varepsilon^2\), so it
is independent of the sign of the additive coupling.  The theorem follows.
\end{proof}

\begin{remark}[scope of the no-go theorems]
The preceding two results close the matching additive class
\[
        \Wcal_{\mathrm{Perelman}}
        +
        \sigma\varepsilon\tau^\alpha K
\]
in the bulk round-collapse limit, and the corresponding matching Planck-wall
additive class at \(K=\KP\).  They do not close arbitrary smooth additive terms
\(H(\tau,K)\); for such a term independent \(H_\tau\) and \(H_K\) contributions
enter the Berger derivative and require a separate admissibility hypothesis or
a separate no-go theorem.  Nor do they rule out genuinely non-additive
noncollapse mechanisms.
\end{remark}





\section{The unmodified Perelman entropy passes the protected Berger cone}

The no-go theorems above concern additive repairs of Perelman's entropy by a
matching \(K\)-potential.  They do not test the unmodified Perelman entropy
itself as an independent noncollapse detector.  On the Berger sector the
unmodified entropy passes the natural two-mechanism test.

Let
\[
        \nu_{\mathrm B}(g)
        =
        \inf_{\tau>0}\,
        \Wcal_{\mathrm{Perelman}}(g,f,\tau)
\]
with \(f\) restricted to the invariant density.  On a homogeneous Berger
metric this density is constant.  In the positive-scalar region
\[
        R=\frac{8-2u}{x}>0
        \quad\Longleftrightarrow\quad
        u<4,
\]
the minimising scale is
\[
        \tau=\frac{3}{2R}.
\]
Using
\[
        \Vol(g_{a,c})=2\pi^2x^{3/2}\sqrt u,
        \qquad
        R=\frac{8-2u}{x},
\]
one obtains, up to an additive universal constant,
\[
        \nu_{\mathrm B}(x,u)
        =
        \log\Vol(g_{a,c})+\frac32\log R+C
        =
        \frac12\log u+\frac32\log(8-2u)+C .
\]

\begin{theorem}[Berger two-mechanism test]
Assume \(\varepsilon>0\) and \(\Phi'(K)>0\).  On the protected
positive-scalar Berger cone
\[
        \frac4{11}\le u<4,
\]
the unmodified Perelman entropy reduction \(\nu_{\mathrm B}\) is monotone along
the curvature-wall Berger flow.  More precisely,
\[
        \frac{d}{dt}\nu_{\mathrm B}
        =
        \frac{32(1-u)^2}{x(8-2u)}
        \left(
        1+\frac{2\varepsilon\Phi'(K)}{x}(11u-4)
        \right)
        \ge0,
\]
with equality only at the round point \(u=1\).
\end{theorem}

\begin{proof}
Differentiating the reduced entropy gives
\[
        \frac{\partial \nu_{\mathrm B}}{\partial u}
        =
        \frac{1}{2u}-\frac{3}{8-2u}
        =
        \frac{4(1-u)}{u(8-2u)}.
\]
The exact Berger anisotropy equation for the wall flow is
\[
        \dot u
        =
        \frac{8u(1-u)}{x}
        \left(
        1+\frac{2\varepsilon\Phi'(K)}{x}(11u-4)
        \right).
\]
Multiplying these two expressions gives the stated derivative formula.  In the
cone \(u\ge4/11\), the bracket is at least one.  Since \(u<4\), the denominator
\(8-2u\) is positive.  Hence \(d\nu_{\mathrm B}/dt\ge0\), and equality occurs
only when \(u=1\).

For the pure quartic wall \(\Phi(K)=K\), this reduces to the fully factored
form
\[
        \frac{d}{dt}\nu_{\mathrm B}
        =
        \frac{16(u-1)^2\bigl[x+2\varepsilon(11u-4)\bigr]}
        {x^2(4-u)} .
\]
For the physical Planck wall
\[
        \Phi(K)=\frac{(K+\delta)^q}{(\KP-K)^p},
\]
one has \(\Phi'(K)>0\) on \(K<\KP\).  Therefore the same monotonicity conclusion
holds throughout the protected positive-scalar cone.
\end{proof}

\begin{corollary}[Berger noncollapse implication under a wall gap]
On the protected positive-scalar Berger cone, suppose in addition that a wall
gap holds:
\[
        K\le \KP-\eta.
\]
Then
\[
        |R|\le \sqrt{3(\KP-\eta)}.
\]
Consequently, on every scale
\[
        0<r\le 3^{-1/4}(\KP-\eta)^{-1/4},
\]
one has \(R\le r^{-2}\).  Since \(\nu_{\mathrm B}(g(t))\ge\nu_{\mathrm B}(g(0))\),
Perelman's static noncollapse implication gives
\[
        \Vol B(x,r)\ge \kappa(\nu_{\mathrm B}(g(0)))\,r^3
\]
on those scales.
\end{corollary}

\begin{remark}[scope]
This is the first positive non-additive Berger result in the present
programme.  It says that the curvature wall and the unmodified Perelman
entropy pass the homogeneous two-mechanism test:
\[
        \text{curvature wall controls high curvature,}
        \qquad
        \nu_{\mathrm B}\text{ controls collapse in the Berger reduction.}
\]
It is not yet a full PDE monotonicity theorem for Perelman's \(\nu\) along the
inhomogeneous DeTurck-gauged wall flow.  The missing analytic statement is a
propagation estimate, or monotonicity formula, for the unmodified Perelman
entropy under the actual fourth-order forced flow.
\end{remark}

\begin{remark}[no PDE envelope identity is used here]
The Berger theorem above is not derived from a formal full-PDE identity for
\(D\nu[g]\).  It uses only the explicit homogeneous minimisation over the
constant density and over \(\tau\), followed by direct differentiation of the
reduced invariant along the verified Berger ODE.

This distinction is important.  A naive formula of the form
\[
        D\nu_g[h]=-\int_M\left\langle
        \Ric+\nabla^2 f-\frac1{2\tau}g,\ h\right\rangle\,d\nu
\]
is not used in this paper.  If \(f\) and \(\tau\) are frozen during the metric
variation, trace and scale terms appear, and a full envelope-theorem derivation
must account for the Euler--Lagrange equations in \(f\) and \(\tau\) with the
chosen normalization.  The inhomogeneous PDE monotonicity problem for the
unmodified entropy is therefore left as a separate analytic problem.
\end{remark}



\section{A PDE programme for the unmodified Perelman entropy}

The Berger theorem proves that the unmodified Perelman entropy is monotone on
the protected positive-scalar Berger cone.  The full inhomogeneous PDE
statement remains open.  This section records the precise analytic and
algebraic problems that must be solved in order to pass from the homogeneous
theorem to a genuine PDE monotonicity or lower-bound propagation theorem.

Let \(M^3\) be closed, and let \(g(t)\) solve the DeTurck-gauged wall flow
\begin{equation}\label{eq:PDE-wall}
\partial_t g = -2\Ric - \varepsilon E_\Phi + \mathcal{L}_W g,
\end{equation}
where
\[
E_\Phi = \Grad_{L^2}\int_M\Phi(|\Rm|^2)\,d\mu,
\qquad
\Phi(K)=\frac{(K+\delta)^q}{(\KP-K)^p}.
\]
Assume on the time interval under consideration that the wall gap holds,
\[
K=|\Rm|^2\le \KP-\eta,
\]
and that the scalar curvature is positive.  Let \((f,\tau)\) be a minimising
pair for Perelman's \(\nu\)-functional at \(g(t)\), and write
\[
d\nu=(4\pi\tau)^{-3/2}e^{-f}\,d\mu,
\qquad
T_{ij}= \Ric_{ij}+\nabla_i\nabla_j f-\frac1{2\tau}g_{ij}.
\]

At a differentiability point of \(\nu\), and for a smooth minimising branch
\((f,\tau)\), the exact first variation is
\[
D\nu_g[h] = -\tau\int_M\langle h,T\rangle\,d\nu .
\]
The DeTurck term contributes nothing, since \(\nu\) is
diffeomorphism-invariant.  Therefore the geometric part of the wall flow gives
the exact identity
\begin{equation}\label{eq:nu-derivative}
\boxed{ \frac{d}{dt}\nu(g(t)) = 2\tau\int_M |T|^2\,d\nu
+ \varepsilon\tau\int_M\langle E_\Phi,T\rangle\,d\nu . }
\end{equation}
Thus the whole PDE monotonicity problem is concentrated in the wall cross term
\[
\mathfrak C_\Phi(g,f,\tau) = \int_M\langle E_\Phi,T\rangle\,d\nu .
\]

The minimiser identities give the weighted divergence relation
\[
\operatorname{div}_f T = 0,\qquad
(\operatorname{div}_f T)_i = \nabla^j T_{ij} - T_{ij}\nabla^j f,
\]
and the averaged trace identity
\[
\int_M\operatorname{tr}T\,d\nu = 0 .
\]
These are weighted identities.  They do not imply that \(T\) is pointwise
trace-free, and they do not imply ordinary divergence-freeness.

The first open task is to compute \(\mathfrak C_\Phi\) exactly from the
unweighted wall tensor
\[
(E_\Phi)_{ij}= \frac12\Phi g_{ij}
-2\Phi'(K)R_i{}^{abc}R_{jabc}
+ 4\nabla^a\nabla^b\bigl(\Phi'(K)R_{iabj}\bigr),
\]
with the convention caveat already stated in the paper.  Since the integral is
weighted by \(d\nu\), the double-divergence term must be integrated by parts
using the weighted adjoint.  This produces terms involving
\[
\nabla f,\qquad \nabla^2 f,\qquad \nabla\Phi'(K),\qquad \nabla^2\Phi'(K),
\]
and curvature commutators.  Under the wall gap, the singular coefficients
\(\Phi'\), \(\Phi''\), and \(\Phi'''\) are bounded by constants depending on
\(\eta\), but the corresponding derivative terms still have to be estimated.

The second open task is algebraic.  One must find a tensorial pinching
condition replacing the Berger condition
\[
u\ge \frac4{11}.
\]
A naive condition such as
\[
|\Ric_0|\le cR,\qquad R>0,
\]
does not automatically reproduce the full protected Berger cone.  The constant
\(c\) must be checked against the Berger model, and the resulting condition
must then be tested for preservation under the fourth-order wall flow.

The desired estimate is either the monotonicity estimate
\[
\mathfrak C_\Phi(g,f,\tau) \ge 0,
\]
or the weaker absorbable bound
\[
\mathfrak C_\Phi(g,f,\tau) \ge -\theta\int_M|T|^2\,d\nu,\qquad
\varepsilon\theta < 2 .
\]
The first estimate would give
\[
\frac{d}{dt}\nu(g(t)) \ge 0 .
\]
The second would still give monotonicity after absorption:
\[
\frac{d}{dt}\nu(g(t)) \ge (2-\varepsilon\theta)\tau\int_M|T|^2\,d\nu \ge 0 .
\]
A still weaker but useful target is an integrated negative-part estimate,
\[
\int_0^T \bigl(\mathfrak C_\Phi(g(t),f_t,\tau_t)\bigr)_- \,dt \le C(T),
\]
which would imply
\[
\nu(g(t)) \ge \nu(g(0)) - C(T).
\]
This lower-bound propagation is sufficient for finite-time noncollapse, even
without strict monotonicity.

The zeroth-order contribution
\[
\frac12\int_M\Phi(K)\operatorname{tr}T\,d\nu
\]
does not vanish automatically, since only \(\int_M\operatorname{tr}T\,d\nu =0\)
is known.  The factor \(\Phi(K)\) prevents immediate cancellation.  This term
must be kept in the cross-term calculation.

Likewise, the curvature term
\[
-2\int_M\Phi'(K)\,T^{ij}R_i{}^{abc}R_{jabc}\,d\nu
\]
has no sign \emph{a priori}.  Its sign or absorbability must be proved from
the same pinching condition that replaces the Berger inequality.

Finally, any PDE monotonicity theorem must be coupled to the analytic theory
of the wall flow: DeTurck parabolicity, short-time existence, preservation of
the wall gap, regularity estimates, and preservation of the required pinching
condition.

Once monotonicity or lower-bound propagation of \(\nu\) is established,
noncollapse follows by Perelman's static theorem.  Indeed, the wall gap
implies
\[
|R|\le \sqrt{3(\KP-\eta)}.
\]
Thus the scalar-curvature hypothesis in the static noncollapse theorem is
satisfied on sufficiently small scales depending only on \(\KP\) and \(\eta\).
A uniform lower bound for \(\nu(g(t))\) then yields
\[
\operatorname{Vol}B(x,r)\ge \kappa\, r^3
\]
on those scales.

The status is therefore:
\[
\boxed{\text{Berger two-mechanism theorem: proved;}}
\]
\[
\boxed{\text{full PDE monotonicity of unmodified \(\nu\): open;}}
\]
\[
\boxed{\text{finite-time noncollapse follows from any propagated lower bound for \(\nu\).}}
\]

Any solution of the full PDE problem must first reproduce the Berger
factorisation in the homogeneous limit.  The Berger theorem is therefore not
merely evidence; it is the consistency test for every proposed inhomogeneous
identity.



\section{Analytic roadmap for the DeTurck-gauged Planck-wall flow}

We record the analytic structure of the curvature-wall flow.  This section is
not a full PDE continuation theorem.  It explains what a future PDE closure
would have to prove.

The ungauged flow is
\[
\partial_t g
=
-2\Ric
-
\varepsilon E_\Phi,
\qquad
E_\Phi
=
\Grad_{L^2}\int_M\Phi(K)\,d\mu,
\qquad
K=|\Rm|^2,
\]
where
\[
\Phi(K)
=
\frac{(K+\delta)^q}{(\KP-K)^p},
\qquad
p>0,\quad q>0,\quad \delta>0 .
\]
With the conventions of this paper, the formal first-variation tensor has
the Streets-type structure
\[
(E_\Phi)_{ij}
=
\frac12\Phi g_{ij}
-
2\Phi'(K)R_i{}^{abc}R_{jabc}
+
4\nabla^a\nabla^b\bigl(\Phi'(K)R_{iabj}\bigr).
\]
The roadmap below uses only its schematic fourth-order part,
\[
4\nabla^a\nabla^b\bigl(\Phi'(K)R_{iabj}\bigr).
\]
The Berger theorems above are independent of a full PDE closure: they use the
explicit homogeneous reduction.  A complete DeTurck theorem would require a
separate convention-by-convention comparison with the standard
\(L^2\)-curvature-flow formula.

After a fourth-order DeTurck gauge, in the sense used for \(L^2\)-curvature
type flows, the metric equation has the schematic leading form
\[
\partial_t g
=
-4\varepsilon\Phi'(K)\Delta^2 g
+
\text{lower-order quasilinear terms},
\]
provided the solution stays in a region
\[
K\le \KP-\eta .
\]
In such a region the coefficients of the equation are smooth.  A full
short-time existence theorem requires the corresponding DeTurck principal
symbol and strong parabolicity verification.  Positivity of \(\Phi'(K)\) is
the expected elliptic sign, but by itself it is not a substitute for the symbol
calculation, since the linearisation also contains terms involving
\(\Phi''(K)\) and the background curvature.

For a metric variation \(v=\partial_tg\), the Riemann tensor satisfies
schematically
\[
\partial_t\Rm
=
\nabla^2v
+
\Rm*v .
\]
Hence the wall part gives
\[
\partial_t\Rm
=
-4\varepsilon\Phi'(K)\Delta^2\Rm
+
\text{terms involving }\nabla\Phi',\nabla^2\Phi',\nabla^j\Rm,
\]
up to the index structure fixed by the Bianchi identities and by the DeTurck
gauge.

Consequently
\[
\partial_tK
=
2\langle\Rm,\partial_t\Rm\rangle
+
\text{metric contraction terms}.
\]
Thus the leading fourth-order contribution is not simply a scalar
maximum-principle term
\[
-4\varepsilon\Phi'(K)\Delta^2K .
\]
Rather, it contains
\[
-8\varepsilon\Phi'(K)
\langle\Rm,\Delta^2\Rm\rangle
\]
plus derivative terms coming from \(\Phi'(K)\).  After integration by parts
these terms are expected to produce coercive fourth-order curvature
dissipation, but pointwise they do not yield a classical maximum principle for
\(K\).

This distinction is essential.  The fourth-order heat operator does not
preserve positivity, and therefore no pointwise comparison argument of the
form
\[
\partial_tK\le F(K)
\]
is available merely from the leading bi-Laplacian structure.

The singular energy nevertheless has a robust variational sign.  Let
\[
\mathcal E(t)=\int_M\Phi(K)\,d\mu .
\]
Along the ungauged flow,
\[
\frac d{dt}\mathcal E
=
\int_M
\langle E_\Phi,-2\Ric-\varepsilon E_\Phi\rangle
\,d\mu.
\]
Hence
\[
\frac d{dt}\mathcal E
=
-\varepsilon\|E_\Phi\|_{L^2}^2
-
2\int_M\langle E_\Phi,\Ric\rangle\,d\mu .
\]
By Cauchy--Schwarz,
\[
-2\int_M\langle E_\Phi,\Ric\rangle\,d\mu
\le
\frac{\varepsilon}{2}\|E_\Phi\|_{L^2}^2
+
\frac{2}{\varepsilon}\|\Ric\|_{L^2}^2 .
\]
Therefore
\[
\boxed{
\frac d{dt}\mathcal E
\le
-\frac{\varepsilon}{2}\|E_\Phi\|_{L^2}^2
+
\frac{2}{\varepsilon}\|\Ric\|_{L^2}^2 .
}
\]
This is the correct energy-dissipation inequality.  It shows that the wall
term dissipates the singular energy up to the Ricci-flow error.  It does not
by itself prove pointwise wall separation.

A precise wall-separation statement requires more than the singular
energy.  It also requires spatial regularity and a lower local volume ratio.
For example, suppose that on a time interval one has
\[
\mathcal E(t)\le E_0,
\qquad
[K(\cdot,t)]_{C^\alpha}\le L,
\]
and, for all sufficiently small \(r\),
\[
\Vol B(x,r)\ge v_0 r^3 .
\]
If
\[
p>\frac3\alpha,
\]
then the singular energy bound forces a positive gap
\[
\sup_MK(\cdot,t)\le \KP-\eta(E_0,L,v_0,p,\alpha).
\]

Indeed, set
\[
m=\KP-\sup_MK .
\]
Choose \(x_0\) with \(K(x_0)\ge \sup_MK-m/2\).  Then
\[
\KP-K(x_0)\le \frac{3m}{2}.
\]
By Hölder control,
\[
K(x)\ge K(x_0)-L\,d(x,x_0)^\alpha .
\]
Taking \(r\asymp (m/L)^{1/\alpha}\), we have
\[
\KP-K(x)\lesssim m
\]
on \(B(x_0,r)\).  Therefore
\[
\mathcal E(t)
\ge
\int_{B(x_0,r)}(\KP-K)^{-p}\,d\mu
\gtrsim
m^{-p}\Vol B(x_0,r)
\gtrsim
v_0 L^{-3/\alpha}m^{3/\alpha-p}.
\]
When \(p>3/\alpha\), this lower bound diverges as \(m\downarrow0\), contradicting
\(\mathcal E(t)\le E_0\).  Hence \(m\) is bounded below by a positive quantity
depending only on \(E_0,L,v_0,p,\alpha\).

Thus the singular wall gives a plausible substitute for a maximum principle
only in combination with spatial regularity and noncollapse-type local volume
control.  The missing analytic step is therefore clear:
\[
\boxed{
\begin{gathered}
\text{energy dissipation }+\text{ regularity/Sobolev control}\\
+\text{ local volume noncollapse }
\Longrightarrow
\text{ pointwise wall separation}.
\end{gathered}
}
\]
This correction is important: the wall energy alone cannot prove noncollapse.
It can help prevent contact with \(K=\KP\) only after enough local volume
control has already been obtained by some independent mechanism.

In the homogeneous Berger reduction no such difficulty appears: all derivative
terms vanish, and the wall reduces to the algebraic inward drift already
computed.  In the full PDE, however, the absence of a fourth-order maximum
principle means that the Berger algebra must be supplemented by DeTurck
parabolicity and a genuine energy-regularity argument.


\section{Remaining noncollapse problem}

The results above have a narrow conclusion.  They identify a local
high-curvature wall, rule out the natural additive Perelman-potential repairs
tested on the Berger sector, and show that the unmodified Perelman entropy
passes the protected positive-scalar Berger cone.  They do not yet produce a
full noncollapse functional for the inhomogeneous PDE.

A complete internal replacement for surgery would require a functional
\(\Ncal\) with the following properties:
\begin{enumerate}[label=(\roman*)]
\item a scale law compatible with the absolute curvature scale
\(\KP=\ell_P^{-4}\);
\item monotonicity along the actual forced DeTurck system;
\item a coercive implication from a lower bound for \(\Ncal\) to local volume
noncollapse under a curvature bound;
\item compatibility with bounded-curvature Hopf-fibre collapse, which is
invisible to the wall \(K\nearrow\KP\).
\end{enumerate}

The preceding results indicate that such an \(\Ncal\), if it exists for the
full PDE, cannot be obtained merely by adding a \(K\)-potential to Perelman's
entropy and using the same \(K\)-gradient to force the metric.  The Berger
calculation points instead to a two-mechanism design: the curvature wall
controls high curvature, while an independent non-additive detector, beginning
with the unmodified Perelman entropy in the homogeneous sector, supplies
noncollapse information.  Other structurally different candidates include
modified reduced volume, Harnack-type quantities for the coupled system,
spacetime action functionals, and heat-kernel detectors.

\begin{openproblem}[PDE propagation of the non-additive detector]
The protected Berger cone is passed by the unmodified Perelman entropy.  The
open problem is to prove an analogue for the full inhomogeneous wall flow:
either a monotonicity formula or a controlled propagation estimate for the
unmodified entropy, sufficient to retain fixed-scale volume noncollapse under
\(K\le r^{-4}<\KP\), and compatible with a DeTurck continuation theorem.
\end{openproblem}

\section{Conclusion}

This paper should be read as a mathematical side branch of the compact \(S^3\)
programme, not as a required physical extension of the operational
Planck-boundary interpretation.  It studies modified Ricci-type flows as
autonomous geometric dynamics and asks which variational curvature barriers
suppress supercritical curvature growth.  The parameter \(\KP=\ell_P^{-4}\)
serves here as a prescribed curvature-wall scale; it is not being promoted to a
new physical forcing law.

A constant cosmological term fails the high-curvature test because its
curvature contribution is quadratic in \(|\Rm|\), while the Ricci reaction is
cubic.  The variational wall
\[
        -\varepsilon\,\grad_{L^2(d\mu)}
        \int_M
        \frac{(|\Rm|^2+\delta)^q}{(\KP-|\Rm|^2)^p}\,d\mu
\]
is the corresponding local high-curvature replacement.  In dimension three it
acts on the full curvature invariant
\[
        K=|\Rm|^2=4|\Ric_0|^2+\frac13R^2,
\]
and its leading conformal drift has the correct sign and degree.

This is not a noncollapse theorem.  The Hopf-fibre degeneration in the Berger
family gives bounded-curvature collapse, so a \(K\)-wall alone cannot replace
surgery.  The exact Berger anisotropy equation shows the same point
dynamically: in the cone \(u\ge4/11\) the round ratio \(u=1\) is attracting,
whereas in the lower Hopf sector the wall can push \(u\) downward when
\[
        \frac{2\varepsilon\Phi'(K)}{x}(4-11u)>1 .
\]
Volume normalization does not alter this anisotropy equation, and the Hopf
limit is an \(S^2\)-quotient collapse rather than a high-curvature limit.
The heat-kernel scale detector \(\mathcal H_r\) diverges along fixed-base Hopf
collapse by the zero-mode lower bound, which identifies a genuinely
non-additive direction for future noncollapse tests.
The natural attempt to turn the wall into an additive Perelman-type
noncollapse entropy fails: \(R+\varepsilon\Phi(K)\) is not monotone even in the
homogeneous Berger sector, and additive matching \(K\)-potential repairs fail
for every sign and every \(\tau\)-weight in the pure quartic test.  The
unmodified Perelman entropy behaves differently: on the protected
positive-scalar Berger cone \(4/11\le u<4\), its reduced derivative along the
wall flow is nonnegative.  This gives a positive non-additive two-mechanism
Berger result, not a full PDE theorem.

The final conclusion is therefore precise:
\[
        \boxed{\text{the curvature wall is a valid local high-curvature regulator;}}
\]
\[
        \boxed{\text{matching additive \(K\)- and Planck-wall potentials do not supply noncollapse;}}
\]
\[
        \boxed{\text{the unmodified Perelman entropy supplies a positive Berger two-mechanism test;}}
\]
\[
        \boxed{\text{a complete surgery replacement still requires PDE propagation of this detector.}}
\]
Thus the contribution of this paper is a mathematical one.  It develops a
curvature-wall mechanism for modified Ricci-type flows and tests it against
Berger collapse, additive entropy repairs, heat-kernel noncollapse indicators,
and the unmodified Perelman entropy.  It should be read as an optional
mathematical refinement, not as a required physical extension of the
operational Planck-boundary interpretation.

Any future candidate, including the unmodified entropy, must pass from the
Berger reduction to an inhomogeneous DeTurck continuation theorem.

\appendix

\section{Algebraic checks}

\subsection{Berger curvature}

With \(x=a^2\), \(z=c^2\), and \(u=z/x\),
\[
        |\Rm|^2
        =
        \frac{4(16a^4-24a^2c^2+11c^4)}{a^8}
        =
        \frac{4(16x^2-24xz+11z^2)}{x^4}
        =
        \frac{4(16-24u+11u^2)}{x^2}.
\]
At \(u=1\),
\[
        |\Rm|^2=\frac{12}{x^2}.
\]
At fibre collapse \(u\to0\) with \(x\) fixed,
\[
        |\Rm|^2\to \frac{64}{x^2}=\frac{64}{a^4}.
\]
At base collapse \(x\to0\) with \(u\) fixed,
\[
        |\Rm|^2\to\infty
\]
because
\[
        11u^2-24u+16>0
\]
for all real \(u\), its discriminant being
\[
        (-24)^2-4\cdot 11\cdot 16=-128<0.
\]

\subsection{Scale exponent}

For a length scale \(L\),
\[
        \int K^q\,d\mu
        \sim
        L^{-4q}L^3
        =
        L^{3-4q}.
\]
The critical exponent is therefore
\[
        3-4q=0
        \quad\Longleftrightarrow\quad
        q=\frac34.
\]

\subsection{Power law}

For \(\Phi(K)=K^q\),
\[
        \Phi'(K)=qK^{q-1},
\]
and hence
\[
        4K\Phi'(K)-3\Phi(K)
        =
        4qK^q-3K^q
        =
        (4q-3)K^q.
\]

\subsection{Barrier law}

For
\[
        \Phi(K)=\frac{(K+\delta)^q}{(\KP-K)^p},
\]
one has
\[
        \frac{\Phi'(K)}{\Phi(K)}
        =
        \frac{q}{K+\delta}
        +
        \frac{p}{\KP-K}.
\]
Thus
\[
        4K\Phi'(K)-3\Phi(K)
        =
        \Phi(K)
        \left(
        \frac{4qK}{K+\delta}
        +
        \frac{4pK}{\KP-K}
        -3
        \right).
\]

\subsection{Sturm check for the Berger falsification}

The polynomial
\[
        c_2(u)=1331u^4-3696u^3+3648u^2-1536u+256
\]
has no positive root.  One Sturm sequence begins with
\[
        c_2,\quad
        c_2'=5324u^3-11088u^2+7296u-1536,
\]
and continues, up to positive scalar multiples, by
\[
        1104u^2-1256u+116,\qquad
        48384-48661u,\qquad
        217220096 .
\]
The sign variation count at \(u=0^+\) equals the sign variation count at
\(u=+\infty\), so the number of positive roots is zero.  Since
\(c_2(0)=256>0\), one has \(c_2(u)>0\) for all \(u>0\).


\section{Reference facts used}

The use of \(L^2\)-curvature-gradient flows as fourth-order degenerate
geometric flows with DeTurck gauge fixing follows the standard theory of the
gradient flow of \(\int_M |\Rm|^2\), especially Streets' DeTurck treatment of
the \(L^2\)-curvature flow.  The heat-kernel coercivity discussion uses the
standard local comparison theory under a Ricci lower bound, in the spirit of
Cheeger--Yau, Li--Yau, and Davies.  The tensorial formulas in the analytic roadmap are
used only at the schematic principal-structure level; a complete PDE theorem
would require a separate convention-by-convention verification against that
standard \(L^2\)-flow formula.  The restriction on second-order natural
divergence-free tensors is the Lovelock-type obstruction: second-order
metric-natural tensors have very limited algebraic freedom, and in low
dimension do not provide a new high-curvature damping mechanism.

\begin{thebibliography}{99}

\bibitem{Hamilton1982}
R.~S. Hamilton,
\emph{Three-manifolds with positive Ricci curvature},
J. Differential Geom. 17 (1982), 255--306.

\bibitem{DeTurck1983}
D.~M. DeTurck,
\emph{Deforming metrics in the direction of their Ricci tensors},
J. Differential Geom. 18 (1983), 157--162.

\bibitem{Perelman2002}
G.~Perelman,
\emph{The entropy formula for the Ricci flow and its geometric applications},
arXiv:math/0211159 (2002).

\bibitem{KleinerLott2008}
B.~Kleiner and J.~Lott,
\emph{Notes on Perelman's papers},
Geom. Topol. 12 (2008), 2587--2855.

\bibitem{ChowLuNi2006}
B.~Chow, P.~Lu, and L.~Ni,
\emph{Hamilton's Ricci Flow},
Graduate Studies in Mathematics 77, AMS, 2006.

\bibitem{BakryEmery1985}
D.~Bakry and M.~Emery,
\emph{Diffusions hypercontractives},
Seminaire de probabilites XIX, Lecture Notes in Math. 1123, Springer, 1985, 177--206.

\bibitem{AngenentKnopf2004}
S.~Angenent and D.~Knopf,
\emph{An example of neckpinching for Ricci flow on \(S^{n+1}\)},
Math. Res. Lett. 11 (2004), 493--518.

\bibitem{Streets2008}
J.~D. Streets,
\emph{The gradient flow of \(\int_M |\Rm|^2\)},
J. Geom. Anal. 18 (2008), 249--271.

\bibitem{Streets2012}
J.~D. Streets,
\emph{Collapsing in the \(L^2\) curvature flow},
Comm. Partial Differential Equations 38 (2013), 985--1014.

\bibitem{CheegerYau1981}
J.~Cheeger and S.-T.~Yau,
\emph{A lower bound for the heat kernel},
Comm. Pure Appl. Math. 34 (1981), 465--480.

\bibitem{LiYau1986}
P.~Li and S.-T.~Yau,
\emph{On the parabolic kernel of the Schrödinger operator},
Acta Math. 156 (1986), 153--201.

\bibitem{Davies1989}
E.~B. Davies,
\emph{Heat Kernels and Spectral Theory},
Cambridge Tracts in Mathematics 92, Cambridge University Press, 1989.

\bibitem{Lauret2018}
E.~A. Lauret,
\emph{The smallest Laplace eigenvalue of homogeneous 3-spheres},
arXiv:1801.04259.


\bibitem{Lovelock1971}
D.~Lovelock,
\emph{The Einstein tensor and its generalizations},
J. Math. Phys. 12 (1971), 498--501.

\bibitem{NavarroNavarro2011}
A.~Navarro and J.~Navarro,
\emph{Lovelock's theorem revisited},
J. Geom. Phys. 61 (2011), 1950--1956.


\bibitem{BergerFalsification2026}
B.~Batenin and A.~Preece,
\emph{The Berger Sector Exactly Falsifies the Naive Noncollapse Candidate
\(\lambda_\Phi=R+\varepsilon\Phi(K)\)},
companion note (2026).

\end{thebibliography}

\end{document}
